% ============================================================================
% main.tex — Buy versus Rent in England: working-paper edition
% ============================================================================
% Compile: xelatex -> biber -> xelatex -> xelatex
% ============================================================================

% !TEX encoding = UTF-8

\documentclass[12pt, a4paper]{article}

% ============================================================================
% wp_header.tex — Standard economics working-paper preamble
% ============================================================================
% Conventions: single column, 12pt, Times-like serif (TeX Gyre Termes),
% 1.5 line spacing, 25mm margins, plain headers, author-year citations,
% no decorative elements. Compile with XeLaTeX + biber.
% ============================================================================

% --- Language ---
\usepackage{polyglossia}
\setdefaultlanguage[variant=british]{english}
\usepackage{csquotes}
\MakeOuterQuote{"}

% --- Geometry ---
\usepackage[a4paper, margin=25mm]{geometry}

% --- Fonts ---
\usepackage{fontspec}
\setmainfont{TeX Gyre Termes}[Ligatures=TeX]

% --- Spacing ---
\usepackage{setspace}
\onehalfspacing
\setlength{\parskip}{0pt}
\setlength{\parindent}{1.5em}

% --- Sectioning: plain bold, standard sizes ---
\usepackage{titlesec}
\titleformat{\section}{\large\bfseries}{\thesection}{0.6em}{}
\titleformat{\subsection}{\normalsize\bfseries}{\thesubsection}{0.6em}{}
\titleformat{\subsubsection}{\normalsize\itshape}{\thesubsubsection}{0.6em}{}
\titlespacing*{\section}{0pt}{1.6em}{0.6em}
\titlespacing*{\subsection}{0pt}{1.2em}{0.4em}
\titlespacing*{\subsubsection}{0pt}{1.0em}{0.3em}

% --- Math ---
\usepackage{amsmath}
\usepackage{amssymb}

% --- Graphics & floats ---
\usepackage{graphicx}
\graphicspath{{figures/}}
\usepackage{float}
\usepackage{placeins}

% --- Tables ---
\usepackage{booktabs}
\usepackage{tabularx}
\usepackage{threeparttable}
\renewcommand{\arraystretch}{1.15}
\AtBeginEnvironment{table}{\small\singlespacing}
\AtBeginEnvironment{threeparttable}{\small}

% --- Captions: plain, bold label ---
\usepackage[font=small, labelfont=bf, format=plain, justification=justified]{caption}

% --- Lists ---
\usepackage{enumitem}
\setlist{nosep, topsep=0.3em, itemsep=0.15em}

% --- Appendix ---
\usepackage[toc,page]{appendix}

% --- Headers: page number only ---
\usepackage{fancyhdr}
\pagestyle{plain}

% --- Footnotes ---
\usepackage[bottom]{footmisc}

% --- Bibliography: author-year ---
\usepackage[
  backend=biber,
  style=authoryear,
  natbib=true,
  maxcitenames=3,
  maxbibnames=10,
  minnames=1,
  giveninits=true,
  uniquename=init,
  url=false,
  doi=false,
  isbn=false,
  dashed=false
]{biblatex}
\setlength{\bibitemsep}{0.5em}
\setlength{\bibhang}{1.5em}
\renewcommand*{\bibfont}{\normalsize\singlespacing}
\setcounter{biburlnumpenalty}{100}
\setcounter{biburlucpenalty}{100}
\setcounter{biburllcpenalty}{100}

% --- Hyperlinks: subdued ---
\usepackage[
  colorlinks=true,
  linkcolor=black,
  citecolor=black,
  urlcolor=blue,
  bookmarks=true
]{hyperref}

% --- Line breaking ---
% A handful of lines (long series codes, bibliography entries, wide tables)
% cannot be broken within the normal tolerance and would otherwise overflow into
% the margin. \emergencystretch lets TeX add interword space on those lines in a
% final pass instead. It affects only lines that would otherwise be overfull.
\setlength{\emergencystretch}{3em}

% --- Utility ---
\newcommand{\gbp}[1]{\mbox{\pounds#1}}

% --- Title page (standard WP style) ---
% \wptitlepage{Title}{Author block}{Date}{Abstract}{JEL}{Keywords}{Thanks footnote}
\newcommand{\wptitlepage}[7]{%
  \begin{titlepage}
    % Author footnote uses a star, per econ convention. NB: \thanks does not
    % work here (its text is only flushed by \maketitle) — use \footnote.
    \renewcommand{\thefootnote}{\fnsymbol{footnote}}%
    \begin{center}
      % The title sits in a field defined by two hairline rules, set one step
      % larger and unbolded so it reads as display type rather than as
      % emphasised body text. Title, rules and abstract share one measure
      % (0.86\textwidth), so the page has a single axis instead of three.
      \vspace*{16mm}
      \rule{0.86\textwidth}{0.4pt}\par
      \vspace{5mm}
      {\LARGE\setstretch{1.08}\selectfont
        \parbox{0.86\textwidth}{\centering #1\par}\par}
      \vspace{5mm}
      \rule{0.86\textwidth}{0.4pt}\par
      \vspace{9mm}
      {\normalsize #2\par}
      \vspace{4mm}
      {\normalsize #3\par}
      \vspace{12mm}
      \begin{minipage}{0.86\textwidth}
        \begin{center}\textbf{Abstract}\end{center}
        \vspace{1mm}
        \noindent {\small\singlespacing #4\par}
        \vspace{5mm}
        {\small\noindent \textit{JEL classification:} #5\par}
        \vspace{1mm}
        {\small\noindent \textit{Keywords:} #6\par}
      \end{minipage}
    \end{center}
    \vfill
    {\footnotesize\noindent #7\par}
  \end{titlepage}
  \setcounter{footnote}{0}%
  \setcounter{page}{2}
}


\addbibresource{references.bib}

\begin{document}

\wptitlepage
  {Buy versus Rent in England:\\ A Cash-Flow-Matched Horse Race\\ for the First-Time Buyer, 2005--2026}
  {Martin Cimprich\footnote{School of Economics and Finance, Queen Mary University of London. Email: \href{mailto:m.cimprich@qmul.ac.uk}{m.cimprich@qmul.ac.uk}. Replication code and data are openly available at \url{https://github.com/Martin-Cimprich/buy-vs-rent-england}. The views expressed are those of the author alone, as are any errors. This paper is independent personal research and does not constitute financial advice.}\\[1mm]
   \normalsize Queen Mary University of London}
  {This version: September 2026}
  {I compare buying a first home with renting and investing, in a cash-flow-matched monthly simulation on official English data from January 2005 to June 2026. The buyer takes a 10\% deposit on a five-year fix, with maintenance and depreciation from the national accounts and first-time-buyer stamp duty; the renter occupies the same dwelling and invests the deposit-equivalent, plus every monthly difference in housing costs, in a global equity index fund, which the comparison assumes they keep up throughout. Price and rent share one population-weighted regional composition, since the official national averages weight regions differently and pairing them biases the result. Over 199 cohorts each holding five years the two strategies finish level, buying ahead in 51\% of cohorts at a median renter/owner wealth ratio of 0.98. Buying wins more often at longer holds, 62\% at seven years and 61\% at ten, but the median renter still finishes within 8--11\% of the owner, \gbp{14{,}454} short over a decade. The mean ratio is 1.22, while the mean gap of \gbp{1{,}744} favours the owner, because renting's proportional wins fell in cohorts where the owner's equity had collapsed, peaking at 4.17 in November 2007. Ownership is competitive only because of the mortgage, since a cash buyer wins 8\% of cohorts, but the same leverage left sixteen northern cohorts entering in 2007--08 with negative equity. At June 2026 a five-year buyer needs 2.3\% appreciation to match equities earning a cautious 5.9\%, and 3.2\% to match the 10.4\% they actually delivered; the dwelling itself appreciated at 3.2\%.}
  {D14, G11, G51, R21, R31}
  {housing tenure choice; buy versus rent; first-time buyers; user cost of housing; household portfolios; leverage; stamp duty}
  {}

% --- Sections ---
% ============================================================================
% 01_introduction.tex
% ============================================================================

\section{Introduction}
\label{sec:introduction}

Buying a first home is among the largest financial decisions most households take, and it carries cultural weight: around 65\% of English households own their home \parencite{ehs2024}, and ownership is widely regarded as financially superior. The public debate on ``buying vs.\ renting'' runs on slogans: renting as ``dead money'' or ``paying your landlord's mortgage'' against weak English house-price returns. Nobody has measured, on monthly data, how a typical first-time buyer fared against an identical renter investing the difference in equities.

This paper runs that race on official English data from January 2005 to June 2026, in a monthly cash-flow-matched simulation after \textcite{beracha2012}. The buyer takes a 10\% deposit and a 30-year term, both at the regulator's median, on a five-year fixed rate at 90\% loan-to-value refixed as the loan amortises, with calibrated maintenance, depreciation, transaction costs and period-accurate stamp duty. The renter occupies the same dwelling at market rent and invests the deposit and purchase costs, plus every monthly difference in outlays, in a global equity index fund (MSCI ACWI, in sterling). I compare the owner's net equity after selling costs with that portfolio.

Three measurement choices shape the results. First, the comparison runs from every monthly entry date, each cohort holding exactly \emph{five years}, roughly a first-time buyer's tenure \parencite{santander_tenure}: 199 cohorts, which isolates \emph{timing}. Second, it is repeated at holding periods from one to ten years. Third, the official England price average reflects the transaction mix while the rent average is dominated by London, so the national price-to-rent ratio (16.9 in June 2026) sits \emph{below every English region} (17.7--22.4) and flatters ownership. I build instead a representative English dwelling, price and rent both population-weighted averages of the nine regions, and run it region by region.

The two strategies were financially comparable, and not only at five years. At five years buying produced greater terminal wealth in 102 of 199 cohorts (51\%), with a median renter/owner wealth ratio of 0.98 and a median gap of \gbp{994} in the buyer's favour. At seven and ten years buying won more often, 62\% and 61\%, but the median renter still finished within 8--11\% of the owner: at ten years the median gap was \gbp{14{,}454} on a dwelling worth \gbp{319{,}488}. Beyond five years buying was ahead more often than not, by margins small enough that renting remained a serious alternative for a household that actually invests the difference every month, which Section~\ref{sec:robustness} shows is where the renter's side of the result comes from.

The distribution behind those medians is asymmetric: buying won more often, renting won bigger. Over five-year holds the mean wealth ratio was 1.22 against a median of 0.98, while the mean gap was \gbp{1{,}744} in the owner's favour. Renting's proportional wins came where the owner's equity had collapsed, so a high ratio sits on a modest sum, while buying won on much larger balance sheets. Renting beat buying twofold in 21 cohorts, threefold in 16 and fourfold in 5, peaking at 4.17 in November 2007: owner \gbp{15{,}769}, renter \gbp{65{,}685}. The ratio's denominator turns negative for some one- and two-year cohorts, so I quote means only from five years out, and never by region, where negative equity leaves a near-zero base.

Timing decides which side of that distribution a household lands on. Renting won every cohort entering in 2006--2009 and 2016; buying won every cohort entering in 2011--2013, 2018 and 2019. The price-to-rent ratio and mortgage rate at purchase jointly account for 56\% of cross-cohort variation, which Section~\ref{sec:predictors} qualifies: the same variables explain 97\% where there is nothing to forecast, so this is largely the user-cost identity; a circular-shift null gives a median $R^2$ of 0.29, leaving $p = 0.034$; and out of sample the model does worse than the historical mean. One regularity survives: renting won all seven cohorts entering above a price-to-rent ratio of 24.4. There is no corresponding floor below which buying was safe: renting also won at the sample's lowest ratio, 17.6. Regional buy-win rates ran from 38\% (North East) to 63\% (West Midlands).

Leverage explains both sides of that record. Global equities returned 10.4\% per year in sterling over the sample against the dwelling's 3.2\%, and the mortgage turns that modest growth into a comparable return on the buyer's own capital; an unlevered cash buyer wins only 8\% of cohorts, and the win rate peaks at 54\% for a 15--20\% deposit against 46\% at 5\%. Leverage also concentrates bad timing: eleven North East and five North West cohorts entering between mid-2007 and mid-2008 ended five years with \emph{negative} net worth after selling costs. The holding-period profile is a feature of this sample rather than of tenure. Buying wins 30\% of one-year cohorts, 51\% at five years and 60--67\% from six to ten; the short end is round-trip transaction costs, roughly thirty per cent of the buyer's initial outlay, which take six years to amortise. Longer holds rest on one crash, one long cheap-credit expansion and one exceptional equity bull market, so I stop at ten years.

Finally, the \textbf{required appreciation rate} (RAR) is the annual house-price growth at which buying exactly breaks even with renting and investing: the user-cost condition of \textcite{poterba1984} and \textcite{himmelberg2005} solved for the capital-gains term. Today's answer turns on what one expects of equities rather than of houses. At June 2026 conditions over five years, the dwelling needs 2.3\% per year to match an equity investor earning a cautious CAPE-implied 5.9\%, but 3.2\% to match one earning the 10.4\% actually delivered. The dwelling's own 3.2\% sits at the top of that range, so the two assumptions bracket the historical outcome, and a single break-even figure would make the case look settled when it is not.

Four contributions follow: the first monthly, cash-flow-matched English backtest with an equity-investing renter, extending the annual, savings-account analysis of \textcite{mostafa2019}, which ends in 2012; the correction to the national price and rent aggregates; the RAR under two openly different views of equity returns; and open code, data and an online implementation.

% ============================================================================
% 02_literature.tex — Related literature
% ============================================================================

\section{Related literature}
\label{sec:literature}

The paper draws on, and speaks to, five strands of work.

\emph{User cost and tenure choice.} The theoretical benchmark for rent-versus-own comparisons is the user cost of owner-occupied housing --- foregone interest on equity, mortgage interest, maintenance and depreciation, less expected capital gains --- formalised by \textcite{poterba1984} on tenure-choice foundations laid by \textcite{henderson1983}, and turned into the standard housing-valuation tool by \textcite{himmelberg2005}. A measurement branch shows, however, that ex-ante user costs diverge persistently from observed rents and are acutely sensitive to imputed appreciation and financing assumptions \parencite{verbrugge2008,diaz2008}, and that credible price-rent comparisons require quality- and composition-matched data \parencite{hill2016,bracke2015}. This paper is the ex-post counterpart to that ex-ante tradition: rather than imputing an equilibrium user cost, I track the realised cash flows of both tenures month by month, which sidesteps the imputation problem; my required appreciation rate is the user-cost equilibrium condition solved for the capital-gains term rather than assuming it; and the composition-consistent national aggregate of Section~\ref{sec:representative} responds directly to the matched-measurement critique --- the official England price and rent aggregates, I show, fail exactly that test.

\emph{Buy-versus-rent horse races.} The closest methodological ancestors are the ex-post backtests for the United States: \textcite{beracha2012}, who run cash-flow-matched buy-versus-rent races over 1978--2009 and find renting frequently competitive; the companion indifference-appreciation exercise of \textcite{berachaseiler2012}, an antecedent of my RAR metric computed under stylised ex-ante assumptions rather than realised paths; and \textcite{rappaport2010} and \textcite{goodman2018} on homeownership as wealth-building, with \textcite{herbert2013} surveying the distributional evidence. For Britain, the only comparable academic exercise is \textcite{mostafa2019}, an annual-frequency comparison of first-time-buyer cohorts over 1975--2011 in which the renter's counterfactual earns building-society deposit rates; buying wins essentially always. \textcite{tabner2016} provides a UK net-present-value framework under assumed parameters, and practitioner analyses exist for Canada \parencite{felix2025}. Relative to this work, the present paper is --- to my knowledge --- the first monthly, cash-flow-matched backtest on UK data, and the first in which the British renter's counterfactual is a diversified equity portfolio; that single change overturns the always-buy conclusion and produces near-parity, while extending the sample through the post-2012 boom, the pandemic, and the 2022--2023 rate shock that earlier UK samples miss.

\emph{Housing returns.} A recent measurement literature quantifies long-run returns to residential property: house prices track incomes rather than compounding like equities over the long run \parencite{knoll2017}; total returns are competitive with equities only once rental yield is included and before costs \parencite{jorda2019rore}; and micro-level studies show that transaction costs, management and idiosyncratic risk substantially erode net returns to individual owners \parencite{chambers2021,eichholtz2021total,giacoletti2021,demers2022}, with large spatial heterogeneity \parencite{amaral2025}. My results are the household-level, England-specific counterpart: with realistic costs the unlevered housing return loses badly to equities --- an all-cash buyer never wins --- and it is mortgage leverage, not the housing return itself, that makes ownership financially competitive.

\emph{Housing in household portfolios.} A structural literature examines how owner-occupied housing shapes household portfolios and lifecycle choices \parencite{flavin2002,cocco2005,yao2005,chetty2017,vestman2019}; \textcite{campbell2006} and \textcite{gomes2021} survey the field. These models emphasise precisely the two forces my backtest measures ex post: the portfolio distortion of a large, levered, indivisible housing position, and the financing channel through which mortgage debt amplifies returns to home equity. The empirical leverage gradient I document --- buying wins two-thirds of cohorts at 95\% LTV and none at all in cash --- is, in effect, a realised-data trace of the mechanism these models study.

\emph{Rent risk, commitment, and the UK market.} Two channels are deliberately switched off in my accounting and frame its interpretation. Ownership hedges rent risk \parencite{sinai2005}, while my renter bears actual market rents throughout; and mortgage amortisation acts as forced saving for present-biased households \parencite{laibson1997,thaler2004,schlafmann2021}, with \textcite{bernstein2024} showing amortisation causally builds net worth and \textcite{sodini2023} identifying substantial benefits of ownership quasi-experimentally --- whereas my renter is endowed with perfect investment discipline. Near-parity for a perfectly disciplined renter therefore implies an advantage for ownership among typical households, which is how I read the result. Finally, the paper connects to UK housing economics --- cycles and regional dynamics \parencite{muellbauer1997,meen1999}, rate-driven valuations \parencite{milesmonro2021}, supply constraints \parencite{hilbervermeulen2016}, and the stamp-duty responses documented by \textcite{besley2014} and \textcite{bestkleven2018}, which motivate my period-accurate SDLT module --- and to the evidence that price-rent ratios forecast housing returns \parencite{caseshiller1989,campbellshiller1988,gallin2008,campbell2009,engstedpedersen2015}: my finding that the entry price-to-rent ratio and mortgage rate explain 74\% of cohort outcomes is the cash-flow-matched analogue of that predictability.

% ============================================================================
% 03_data.tex — Data
% ============================================================================

\section{Data}
\label{sec:data}

The analysis draws on official statistics for all housing inputs and on index-provider data for financial returns. All series are monthly, in nominal pounds sterling, and cover January 2005 to June 2026 (258 months). I focus on England and its nine regions: England is the only part of the United Kingdom for which official monthly rent statistics reach back to 2005.

The sample is bounded at both ends by data rather than by choice, and the two ends bind for different reasons. The start is set by rents: no regional private-rent series for England exists at monthly or quarterly frequency before January 2005, in any source, official or commercial. The end is set by prices: the UK House Price Index runs about two months in arrears, so June 2026 is the latest complete vintage even though rents are published to July. Appendix~\ref{app:data} provides construction details; Table~\ref{tab:sources} summarises the sources.

\begin{table}[H]
  \centering
  \caption{Data sources}
  \label{tab:sources}
  \begin{threeparttable}
  \begin{tabularx}{\textwidth}{@{}lXl@{}}
    \toprule
    \textbf{Series} & \textbf{Source} & \textbf{Coverage used} \\
    \midrule
    House prices & UK House Price Index, HM Land Registry (all dwellings, England and regions) & 2005:1--2026:6 \\
    Private rents & ONS Price Index of Private Rents (PIPR) and its historical series & 2005:1--2026:6 \\
    Mortgage rate & Bank of England quoted rates: five-year fix at 90\% loan-to-value, constructed (Section~\ref{sec:mortgage_rate}) & 2005:1--2026:6 \\
    Equity returns & MSCI ACWI net total return, GBP, less 0.12\% p.a.\ fund fee & 2005:1--2026:6 \\
    Gilt yield & 10-year UK government bond yield (OECD, via FRED) & 2005:1--2026:5 \\
    Margin rate & Bank of England Bank Rate plus 1.50\% (renter-leverage robustness) & 2005:1--2026:6 \\
    Stamp duty & HMRC SDLT schedules, first-time-buyer rules by month & 2003:12--2026:6 \\
    \bottomrule
  \end{tabularx}
  \end{threeparttable}
\end{table}

\subsection{House prices and rents}

House prices come from the UK House Price Index (UK HPI), the official transaction-based index constructed from HM Land Registry sales records by hedonic regression with geometric averaging \parencite{ukhpi}. I use the average price of all dwellings, monthly, for England and each region.

Rents come from the ONS Price Index of Private Rents (PIPR), the official measure of private rents based on Valuation Office Agency lettings data \parencite{ons_pipr,ons_bulletin2026}. Two features make PIPR the appropriate counterpart to the UK HPI. First, it publishes average rent \emph{levels} in pounds per calendar month, not only an index. Second, it is a \emph{stock} measure covering new and existing tenancies, matching the position of a long-term renter whose rent adjusts with the market rather than only at new lets. The live PIPR series begins in January 2015; earlier years use the ONS historical series, chain-linked to PIPR's predecessor index back to January 2005 for England and its regions \parencite{ons_pipr_hist}. The two sources agree exactly over their 122-month overlap, so the splice introduces no discontinuity.

\subsection{A composition-consistent representative dwelling}
\label{sec:representative}

Moving from regions to a single national figure raises a significant measurement issue. The headline UK HPI England price and the headline PIPR England rent each average the country, but they weight the regions differently: the price is a transaction-mix average, close to a simple average of the nine regions, whereas the rent is a tenancy-mix average dominated by South East and London's large private rental sector. In June 2026 the national average rent (\gbp{1{,}446}) exceeded the average rent of every region except London, while the national average price (\gbp{293{,}262}) sat near the middle of the regional range. The ratio of the two, a price-to-rent ratio of 16.9, consequently falls \emph{below the ratio of every individual region} (17.7 to 22.4). Any rent-versus-own comparison built directly on this pair prices a mid-England home against a more expensive London-weighted rent, overstating the cost of renting and biasing the comparison towards ownership.

I therefore construct the price and rent of a \emph{representative English dwelling} as the same population-weighted average of the nine regions:
\begin{equation}
  P^{\text{rep}}_t = \sum_{i=1}^{9} w_i\, P_{i,t}, \qquad
  R^{\text{rep}}_t = \sum_{i=1}^{9} w_i\, R_{i,t}, \qquad
  w_i = \frac{\text{pop}_i}{\sum_j \text{pop}_j},
  \label{eq:representative}
\end{equation}
with fixed weights from mid-2023 ONS regional population estimates. Both aggregates then share one regional composition, and the resulting price-to-rent ratio (20.7 in June 2026) lies within the regional range. Over the sample the representative dwelling appreciated from \gbp{162{,}312} to \gbp{319{,}488} (3.2\% per year), falling 18\% peak-to-trough during 2007--2009, while its rent rose from \gbp{692} to \gbp{1{,}286} per month (2.9\% per year). The regional analyses in Sections~\ref{sec:results} and~\ref{sec:rar} use each region's own price and rent and are unaffected by aggregation.

\subsection{Mortgage rates}
\label{sec:mortgage_rate}

The rate a first-time buyer pays is not the rate the average borrower pays. English first-time buyers overwhelmingly fix: of new first-time-buyer lending in the first half of 2025, 99.1\% was fixed-rate and 0.018\% originated on a reversion rate, and among the fixed loans the three-to-five-year band was the modal choice in six of the seven half-years the Financial Conduct Authority publishes \parencite{fca_psd}. The baseline borrower therefore takes a \textbf{five-year fix at 90\% loan-to-value}, refixed at the prevailing five-year rate every sixty months.

No single Bank of England series measures that: the five-year 90\% loan-to-value series begins in February 2019. Earlier months are constructed by interpolating in loan-to-value between the published five-year series at 75\% and 95\%, with a sixty-one-month gap during 2008--2013, when fewer than three lenders offered a 95\% five-year fix, filled from the two-year 90\% series. Appendix~\ref{app:data} gives the construction and its validation. The series averages 4.74\% and ranges from 2.33\% to 6.99\%.

Two features of it matter for the results. It sits 1.19 percentage points above the Bank's effective rate on all new advances (CFMBJ95). That series is volume-weighted across every loan-to-value tier and fixation length, so recasting it monthly removes the fixation cliff through which a British first-time buyer is actually exposed to rate moves; Section~\ref{sec:robustness} reports it as a bounding case, and the difference is large. And because lenders price in bands and a borrower moves between them, the rate is repriced at each refix by the borrower's \emph{current} loan-to-value. A buyer entering at 90\% on a thirty-year term is at 74.8\% by their first refix, and charging them the 90\% premium then, a premium that peaked at 2.2 percentage points over 2010--2014, would overstate the cost of owning by several percentage points in exactly the period that decides the comparison.

\subsection{Financial returns}

The renter's portfolio tracks the MSCI ACWI index, covering large- and mid-capitalisation equities in 47 developed and emerging markets \parencite{msci_acwi}. I use end-of-month index levels computed by MSCI directly in pounds sterling in the net-total-return variant (dividends reinvested after withholding taxes), less an ongoing charge of 0.12\% per year representing a low-cost global index fund available to UK retail investors. Measuring the index in sterling embeds the currency exposure an unhedged UK investor bears. Over the sample the net-of-fee series returned 10.4\% per year, a factor of 8.38. As a home-bias robustness check I also use a FTSE~100 total-return series, which returned 7.1\% per year, a factor of 4.38; the choice between the two turns out to be one of the most consequential in the paper, and Section~\ref{sec:robustness} reports it in full. The 10-year gilt yield (OECD, via FRED) serves as a conservative alternative opportunity cost in Section~\ref{sec:rar}; it ranged from 5.4\% (June 2007) to 0.2\% (July 2020) and stood at 4.9\% in May 2026, its latest observation \parencite{fred_uk_gilts}.

Figure~\ref{fig:market_overview} shows the model's housing inputs together with the price-to-rent ratio they imply; Figure~\ref{fig:stock_returns} compares cumulative asset performance. The gap the buyer must overcome is evident: over twenty-one years the global equity portfolio multiplied roughly eightfold, UK large-cap equities fourfold, and the representative English dwelling, before the owner's running costs, twofold.

\begin{figure}[H]
  \centering
  \includegraphics[width=\textwidth]{fig1_market}
  \caption{Housing-market inputs for the representative English dwelling, 2005--2026. Panel A: average house price. Panel B: average private rent. Panel C: price-to-rent ratio (price divided by annual rent); the dashed line marks the sample mean of 21.9. Panel D: Bank of England effective interest rate on new mortgage advances to households. Sources: HM Land Registry UK HPI; ONS PIPR; Bank of England.}
  \label{fig:market_overview}
\end{figure}

\begin{figure}[H]
  \centering
  \includegraphics[width=\textwidth]{fig2_assets}
  \caption{Cumulative performance of global equities (MSCI ACWI, net total return in GBP, less a 0.12\% fund fee), UK equities (FTSE 100 total return) and representative English house prices, January 2005 = 100. Sources: MSCI; Morningstar; HM Land Registry.}
  \label{fig:stock_returns}
\end{figure}

% ============================================================================
% 04_framework.tex — Empirical framework
% ============================================================================

\section{Empirical framework}
\label{sec:framework}

The simulation compares two housing tenure strategies on a cash-flow-matched basis: both start with identical capital, and their cash flows are equalised every month. Whenever the renter pays less for housing in a given month, the difference is invested; whenever the renter pays more, the difference is withdrawn from the portfolio. The design follows \textcite{beracha2012}, with transaction costs and taxes calibrated to the English first-time buyer.

\subsection{The buyer}

The buyer purchases the property at the market price prevailing in the entry month. In the baseline scenario, they put down a deposit of 10\% of the price (loan-to-value of 90\%), pay purchase costs of 1.2\% of the price (conveyancing, survey, mortgage and valuation fees), and stamp duty under the first-time-buyer schedule in force in that month. The mortgage amortises over 30 years.

The baseline rate is a five-year fix, because 99\% of first-time-buyer lending is fixed-rate and the five-year band is the modal choice. The payment is set at entry from the quoted five-year rate at the buyer's loan-to-value, held constant for 60 months, then reset at the rate prevailing at the refix:
\begin{equation}
  r_{m,t} = (1 + r_{t})^{1/12} - 1, \qquad
  \text{Payment}_t = \frac{B_t\, r_{m,t}}{1 - (1 + r_{m,t})^{-(T-t)}},
  \label{eq:payment}
\end{equation}
where $r_t$ is the annual rate in force, $B_t$ the outstanding balance and $T-t$ the remaining months. At each refix the rate is taken from the band matching the loan-to-value the buyer has amortised down to, so a borrower who started at 90\% is not charged the high-loan-to-value premium indefinitely. Appendix~\ref{app:data} sets out the rate construction. A floating rate and a two-year fix are examined in Section~\ref{sec:robustness}, where the fixation design turns out to move the answer more than any parameter except maintenance.

Each month the owner also pays maintenance and depreciation of 1.5\% per year of the property's current market value,
\begin{equation}
  M_t = \frac{0.015}{12}\, V_t,
  \label{eq:maintenance}
\end{equation}
an all-in charge covering routine maintenance, major works and replacements, buildings insurance, and physical depreciation not offset by spending. It is the single most consequential assumption in the paper, so it is built from the national accounts rather than from convention.

The charge has two components. \emph{Depreciation} is measured by ONS consumption of fixed capital on dwellings, divided by the market value of the dwelling stock including the land beneath it: 1.20\% per year on average over 1995--2024, 1.28\% in 2024, with a range of 0.95\% to 1.47\% \parencite{ons_cfc,ons_balance_sheet}. \emph{Cash maintenance and buildings insurance} add roughly 0.35\%, from ONS Family Spending \parencite{ons_family_spending}. Together, 1.55\%, which the baseline rounds to 1.5\%.

The land share is what makes this figure specific to the United Kingdom, and it is why figures imported from the American literature do not transfer. Depreciation applies to structures, not to the ground they stand on, and in the UK land is 69.9\% of the market value of household residential property. A rate estimated on structure value, or on total value in a market where land is half of it, therefore overstates the charge here by a wide margin once applied to a British price. On a UK basis, ONS implies depreciation of 4.24\% of \emph{structure} value, which is 1.28\% of market value at a 30.1\% structure share.

Section~\ref{sec:robustness} reports 1.2\% to 1.8\% as the range UK evidence supports, the floor being depreciation alone, for a buyer who does no discretionary upkeep, and the ceiling being the peak observed depreciation year plus cash maintenance. Values of 1\% and 2--3\% are also reported, but as bounds beyond the evidence rather than as equally supported alternatives. The property value itself evolves with the relevant house-price index, $V_{t+1} = V_t (1 + g_t)$.

The buyer's terminal wealth after a holding period ending at $T^{*}$ is net housing equity after selling costs of 1.8\% (estate-agent commission, conveyancing, energy certificate):
\begin{equation}
  W^{\text{own}} = V_{T^{*}} (1 - 0.018) - B_{T^{*}}.
  \label{eq:nw_buy}
\end{equation}

\subsection{The renter}

The renter occupies an equivalent dwelling at the market rent prevailing each month and invests all spare cash. The initial portfolio equals the buyer's total initial outlay (deposit plus purchase costs plus stamp duty), and each month the difference between the owner's outflow (mortgage payment plus maintenance) and the rent is added to, or withdrawn from, the portfolio, which earns the monthly return on the global equity index (net of withholding taxes and fund fees). The renter's terminal wealth is the portfolio value, $W^{\text{rent}} = P_{T^{*}}$. 

\subsection{Outcome measure and analyses}

Outcomes are summarised by the renter/owner wealth ratio
\begin{equation}
  R = \frac{W^{\text{rent}}}{W^{\text{own}}},
  \label{eq:ratio}
\end{equation}
following \textcite{felix2025}; $R > 1$ means renting-and-investing produced greater terminal wealth. The ratio is invariant to the initial capital scale and comparable across cohorts and regions.

Four analyses are run: (i) 199 entry cohorts, one for every month from January 2005 to July 2021, each holding for exactly five years, which is the primary result; (ii) the same exercise repeated at every holding period from one to ten years; (iii) both region by region, aggregated with population weights; and (iv) the forward-looking required appreciation rate of Section~\ref{sec:rar}, at end-of-sample conditions.

The horizon panel stops at ten years because the sample cannot support more. A fifteen-year hold can only be evaluated for cohorts entering in the first seventy-nine months and a twenty-year hold only for the first nineteen, all of them buying into the pre-crisis peak, so those figures would describe one entry window rather than the effect of holding longer. Ten years still rests on 139 cohorts spanning eleven and a half years of entry dates.

Cohorts start every month rather than every quarter, which removes an arbitrary grid at no cost. It should not be mistaken for extra evidence: the underlying sample is the same 258 months either way, and it contains only 4.3 non-overlapping five-year windows however often the simulation is restarted. Cohort counts in this paper are counts of overlapping simulations, not of independent observations, and Section~\ref{sec:predictors} takes that seriously where it matters.

\subsection{Stamp duty}

Stamp Duty Land Tax is modelled period-accurately: each simulated purchase pays the tax implied by the England first-time-buyer schedule in force in that month, including the 2008--2009 holiday, the 2010--2012 first-time-buyer relief, the 2014 slab-to-slice reform, the permanent relief from November 2017, the 2020--2021 holiday, and the 2022--2025 threshold changes \parencite{govuk_sdlt,hoc_sdlt}. Appendix~\ref{app:sdlt} tabulates the schedules. For the representative dwelling, a first-time buyer paid SDLT in 143 of 258 purchase months (55\%), averaging \gbp{1{,}825} (0.88\% of the price) when payable and never exceeding \gbp{3{,}144}. Stamp duty is therefore a second-order factor nationally. It is not second-order in London, where the average dwelling now sits above the \gbp{500{,}000} ceiling at which first-time-buyer relief is withdrawn entirely, and the bill jumps to \gbp{17{,}694} (Section~\ref{sec:rar}).

\subsection{Scope and abstractions}

The model excludes several factors by design. Income and capital-gains taxes are omitted on both sides: a first-time buyer's main residence is exempt from capital-gains tax, and the renter's flows fit within annual ISA allowances for most of the sample, so the symmetric no-tax treatment is a reasonable approximation for the typical household (though not for very large portfolios held outside tax wrappers). Council tax and utilities are paid by the occupant under both tenure options and net out. In the baseline the renter never borrows to invest, so leverage is available only through the mortgage; Section~\ref{sec:robustness} relaxes that and lets the renter borrow on margin. Non-financial considerations, among them security of tenure, freedom to modify, flexibility to move and the insurance value of ownership against local rent risk \parencite{sinai2005}, are outside the accounting and are discussed in Section~\ref{sec:discussion}. Table~\ref{tab:assumptions} collects the baseline parameters.

Two things about that table are easy to mistake for claims about the world.

First, \textbf{every value in it is a baseline and not a prediction about any individual household}. Each has been calibrated to the best available evidence on what an English first-time buyer typically does: the deposit share and term to Financial Conduct Authority loan-level data, the fixation design to the regulator's published breakdown of fixed-rate lending, the transaction costs to first-time-buyer fee surveys, the maintenance charge to national-accounts depreciation. But a typical value is not a universal one. Deposits range from 5\% to 25\%, terms from 25 to 40 years, and a buyer's own figures may sit anywhere in those ranges. Section~\ref{sec:robustness} reports how the result moves as each parameter varies across its defensible range.

Second, the parameters are not equally consequential. The maintenance charge, the fixation design and the renter's portfolio move the answer a great deal; the transaction costs and the mortgage term barely move it. Section~\ref{sec:robustness} ranks them.

\begin{table}[H]
  \centering
  \caption{Baseline parameters (first-time buyer)}
  \label{tab:assumptions}
  \begin{threeparttable}
    \begin{tabularx}{\textwidth}{@{}lXr@{}}
      \toprule
      \textbf{Parameter} & \textbf{Description} & \textbf{Value} \\
      \midrule
      Property              & Representative English dwelling (Section~\ref{sec:representative}) & --- \\
      Holding period        & Typical tenure in a first home                     & 5 years \\
      Deposit               & Share of purchase price (90\% loan-to-value)       & 10\% \\
      Mortgage term         & Amortisation period                                & 30 years \\
      Rate type             & Five-year fix, refixed at the prevailing five-year rate, repriced by current loan-to-value & --- \\
      Purchase costs        & Conveyancing, survey, mortgage product fee         & 1.2\% \\
      Stamp duty            & England first-time-buyer schedule, by month        & rule-based \\
      Selling costs         & Estate agent, conveyancing, EPC                    & 1.8\% \\
      Maintenance \& depreciation & Share of current market value, per year      & 1.5\% \\
      Equity index          & MSCI ACWI net total return, GBP, less 0.12\% fee   & --- \\
      Renter's leverage     & Unlevered portfolio                                & 1.0$\times$ \\
      Rent                  & Market rent, updated monthly                       & --- \\
      Currency              & Nominal GBP throughout                             & --- \\
      \bottomrule
    \end{tabularx}
  \end{threeparttable}
\end{table}

% ============================================================================
% 05_results.tex — Results
% ============================================================================

\section{Results}
\label{sec:results}

\subsection{One long race, as an illustration}
\label{sec:single}

A buyer and a renter both start in January 2005 and are evaluated in June 2026. The buyer purchases the representative dwelling for \gbp{162{,}312} with a deposit of \gbp{16{,}231} and purchase costs of \gbp{3{,}571} including \gbp{1{,}623} of stamp duty; the renter starts with the identical \gbp{19{,}802} in the equity index and occupies the same dwelling at \gbp{692} per month. This is an illustration, not a result: a twenty-one-year hold is far longer than a first-time buyer keeps a first home, and one entry date says little about a decision that turns on timing. Figure~\ref{fig:single_wealth} shows the two wealth paths.

The renter finishes 44\% ahead, with \gbp{367{,}821} against the owner's \gbp{254{,}658} (a \gbp{319{,}488} dwelling less \gbp{59{,}079} of outstanding mortgage and \gbp{5{,}751} of selling costs). Figure~\ref{fig:cashflows} shows the matched flows. The owner's outlay exceeded the rent in 188 of 258 months, so the renter added to the portfolio in most of them: the mortgage payment starts at \gbp{801} against rent of \gbp{692}, and maintenance grows from \gbp{203} to \gbp{399} per month as the property appreciates. Because the rate refixes every five years rather than floating, the payment moves in steps and peaks at \gbp{820} instead of tracking the 2009 collapse downwards; by June 2026 it is \gbp{687} against rent of \gbp{1{,}286}.

A leveraged claim on an asset appreciating at 3.2\% per year returned \gbp{254{,}658} on \gbp{19{,}802} of capital, against \gbp{367{,}821} from a portfolio compounding at 10.4\%. Unlevered, the same dwelling wins 8\% of five-year cohorts (Section~\ref{sec:robustness}).

\begin{figure}[H]
  \centering
  \includegraphics[width=\textwidth]{fig3_wealth}
  \caption{Terminal wealth of the two strategies, January 2005 entry cohort held to June 2026. The buyer's wealth is the market value of the dwelling less the outstanding mortgage; the renter's is the value of the investment portfolio. Selling costs are applied at the terminal date only.}
  \label{fig:single_wealth}
\end{figure}

\begin{figure}[H]
  \centering
  \includegraphics[width=\textwidth]{fig7_cashflows}
  \caption{Monthly housing outlays of the two strategies, January 2005 entry cohort. The owner pays the mortgage payment plus maintenance; the renter pays the market rent. Shading marks the gap that is invested (owner pays more) or withdrawn (renter pays more) each month under cash-flow matching.}
  \label{fig:cashflows}
\end{figure}

\subsection{Entry cohorts over a five-year hold}
\label{sec:rolling}

The single 2005 start is one draw from history. To isolate timing, I repeat the race for every entry month, each cohort holding for exactly five years.

Five years is chosen because it is roughly how long a first-time buyer keeps their first home: Santander put the average at 4.5 years \parencite{santander_tenure}. Fixing the holding period also makes the accounting coherent. Every cohort faces the same round trip (stamp duty and 1.2\% purchase costs on the way in, 1.8\% selling costs on the way out), and the selling costs are charged at the moment the buyer actually incurs them, when they move on. The alternative, running every cohort forward to the end of the sample, produces cohorts of wildly unequal length; Section~\ref{sec:robustness} reports it.

This gives 199 cohorts entering between January 2005 and July 2021. Buying produced greater wealth in 102 of them (51\%) and renting in 97 (49\%). The median wealth ratio is $R = 0.980$ and the median wealth gap \gbp{994} in the buyer's favour. On the central case the two strategies finish level. The cohorts overlap heavily, so these are counts of simulations rather than of independent histories: 258 months of data contain only 4.3 non-overlapping five-year windows, a point Section~\ref{sec:predictors} returns to.

\begin{figure}[H]
  \centering
  \includegraphics[width=\textwidth]{fig4_cohorts}
  \caption{Cohort outcomes. Panel A: renter/owner wealth ratio $R$ by entry month, 199 cohorts each holding for five years. Bars extend from unity; values above one indicate that renting-and-investing produced greater terminal wealth. Panel B: the share of cohorts won by buying at each holding period from one to ten years (left axis), with the median wealth ratio at each horizon on the right axis; monthly entry months at every horizon. Baseline parameters of Table~\ref{tab:assumptions}.}
  \label{fig:rolling}
\end{figure}

Around that parity the dispersion is wide, systematic and asymmetric (Figure~\ref{fig:rolling}, Panel A), and the mean and the median of the 199 cohorts point in opposite directions. The mean wealth ratio is 1.220 against a median of 0.980, so the average cohort left the renter a fifth ahead in proportional terms. The mean wealth gap is \gbp{1{,}744} in the owner's favour against a median of \gbp{994}, so the average cohort left the owner ahead in pounds. Both figures are correct, and the reason they disagree is where each strategy won. Renting's large proportional wins came in cohorts where the owner's equity had collapsed to a small number, and a fourfold ratio over a small sum is a modest amount of money; buying's wins came on much larger balance sheets, where a ratio close to one can still stand for tens of thousands of pounds. The ratio measures how many times over one strategy beat the other and the pound gap measures by how much, so I quote both. The mean ratio on its own would overstate renting's case.

The mean is well defined at this horizon because the owner's terminal net worth never fell below \gbp{15{,}769} across the 199 cohorts, so the denominator of $R$ stays clear of zero. That does not hold region by region, and Section~\ref{sec:regional} reports medians only for the reason set out there.

The tails are lopsided. Renting finished more than twice as wealthy in 21 cohorts, more than three times in 16 of those and more than four times in five; the five all entered between August and December 2007, and the sixteen between April 2007 and July 2008. The peak is the November 2007 cohort at $R = 4.165$, where the owner finished five years with \gbp{15{,}769} against the renter's \gbp{65{,}685}, a gap of \gbp{49{,}916}; that owner supplies the sample minimum quoted above. Buying's best cohort is May 2013 at $R = 0.639$, with the owner on \gbp{100{,}336} against the renter's \gbp{64{,}123}, a gap of \gbp{36{,}214}. Buying won more cohorts and renting won the larger ones, and the largest of renting's wins are proportional wins over small absolute sums. A 10\% deposit is what makes ownership competitive at all against equities compounding at 10.4\% a year, and it is also what turns one mistimed entry into the near-total loss of that deposit.

Timing sorts the cohorts by entry year, and not smoothly. Buying won 17\% of cohorts entering in 2005 and none at all of those entering in 2006, 2007, 2008 or 2009. Its share then rises through 33\% in 2010 to every cohort entering in 2011, 2012 and 2013, and 83\% of those entering in 2014. The run breaks twice: 2015 split evenly, and renting won every cohort entering in 2016, whose five years ran through the equity gains of 2019--2021. Buying then took 92\% of 2017 entries, every cohort entering in 2018 and 2019, 67\% in 2020, and one of the seven 2021 cohorts the five-year panel can evaluate, whose holds contain the 2022--2023 rate shock. Stamp duty is a minor factor: 120 of the 199 cohorts paid any, never exceeding 1.2\% of the purchase price.

\subsubsection*{Other holding periods}

Panel B of Figure~\ref{fig:rolling} repeats the exercise at every holding period from one to ten years. Buying wins 30\% of one-year cohorts, 47\% at four years, 51\% at five and 60--67\% from six to ten, so past five years it is ahead more often than not. The median moves far less than that share does: the median ratio is 0.918 at seven years and 0.891 at ten, leaving the median renter within 8\% and 11\% of the owner, and the median ten-year gap is \gbp{14{,}454} on a dwelling worth \gbp{319{,}488}. The mean ratio is 1.031 at seven years and 0.963 at ten, and the seven-year figure carries the same caveat as the five-year one: the median gap there is \gbp{9{,}272} in the owner's favour. I do not quote it below five years. The minimum owner net worth is negative at one and two years ($-$\gbp{12{,}346} and $-$\gbp{8{,}822}), so the denominator of $R$ changes sign, which is why the mean falls below the median at two years and spikes to 1.471 at three. Panel B therefore plots the median ratio.

The short end of the profile is transaction costs. A buyer pays stamp duty plus 1.2\% on entry and 1.8\% on exit, roughly three per cent of the purchase price in round-trip costs against a 10\% deposit, so the transaction itself consumes about thirty per cent of the initial outlay. Over one year there is almost no appreciation or amortisation to set against that, and the buyer loses on arithmetic alone. Each additional year spreads the same fixed cost more thinly, which is why the curve rises steeply to about year six.

This profile is a property of the sample rather than a general fact about tenure. Twenty-one years of English data contain one house-price crash, one long expansion of cheap credit and one exceptional run in global equities, and the shape of the curve depends on where those three episodes fall inside a given holding period. Ten years is also as far as the data reach at a usable cohort count, 139 against 199 at five years, so I report nothing longer.

\subsection{Regional results}
\label{sec:regional}

Running the race region by region, each with its own price \emph{and} rent so that each is internally consistent, validates the representative-dwelling aggregate. Table~\ref{tab:regional} reports five-year cohort outcomes and the forward break-even for each region.

A note on the statistic. $R$ is a ratio, and with a 10\% deposit the buyer's terminal net worth can cross zero, at which point the ratio explodes: in the North East it ranges from $-82$ to $+9{,}567$ across the 199 cohorts, so its \emph{mean} is not a usable summary. That condition holds nationally and fails here. For the representative dwelling the owner's net worth never fell below \gbp{15{,}769}, which is why Section~\ref{sec:rolling} can quote a mean ratio and this table cannot. I therefore report the median ratio and, alongside it, the median wealth gap in pounds, which cannot behave that way. The share of cohorts won by buying is unaffected by the problem.

\begin{table}[H]
  \centering
  \caption{Regional outcomes over five-year holds, 2005--2026}
  \label{tab:regional}
  \begin{threeparttable}
  \small
  \setlength{\tabcolsep}{4.2pt}
  \begin{tabular}{@{}lrrrrrrr@{}}
    \toprule
    & \textbf{Price} & & \textbf{Buy-} & \textbf{Median} & \textbf{Median} & \textbf{5y break-} & \textbf{Realised} \\
    \textbf{Region} & \textbf{(2026:6)} & \textbf{P/R} & \textbf{win \%} & \textbf{$R$} & \textbf{gap} & \textbf{even} & \textbf{appr.} \\
    \midrule
    North East               & \gbp{165{,}550} & 17.7 & 38 & 1.11 & $+$\gbp{6{,}309}  & 1.17\% & 2.13\% \\
    North West               & \gbp{219{,}922} & 19.1 & 55 & 0.77 & $-$\gbp{7{,}127}  & 1.69\% & 3.26\% \\
    London                   & \gbp{553{,}870} & 20.1 & 43 & 1.12 & $+$\gbp{16{,}223} & 2.82\% & 3.81\% \\
    Yorkshire and The Humber & \gbp{207{,}948} & 20.1 & 52 & 0.93 & $-$\gbp{2{,}923}  & 2.02\% & 2.93\% \\
    West Midlands            & \gbp{234{,}731} & 20.1 & 63 & 0.81 & $-$\gbp{12{,}249} & 2.03\% & 3.00\% \\
    South West               & \gbp{304{,}562} & 20.5 & 48 & 1.01 & $+$\gbp{701}      & 2.16\% & 2.77\% \\
    East Midlands            & \gbp{240{,}457} & 21.8 & 56 & 0.95 & $-$\gbp{3{,}327}  & 2.49\% & 2.90\% \\
    East of England          & \gbp{338{,}707} & 22.0 & 49 & 1.01 & $+$\gbp{1{,}089}  & 2.69\% & 3.17\% \\
    South East               & \gbp{380{,}380} & 22.4 & 42 & 1.08 & $+$\gbp{6{,}575}  & 2.90\% & 3.12\% \\
    \midrule
    England (pop.-weighted)  & --              & 20.6 & 49 & 0.98 & $+$\gbp{1{,}263}  & 2.34\% & 3.12\% \\
    \bottomrule
  \end{tabular}
  \begin{tablenotes}
    \small
    \item \textit{Notes:} 199 cohorts per region, each entering in a distinct month and holding for five years. Buy-win \% is the share in which buying produced greater terminal wealth. Median $R$ is the median renter/owner wealth ratio ($R>1$: renting won); median gap is the median of renter minus owner terminal wealth, so a negative figure favours the buyer. 5y break-even is the forward required appreciation against equities at the cautious 5.9\% assumption (Section~\ref{sec:rar}). Realised appr.\ is the region's own annualised house-price growth over 2005--2026. P/R is the June 2026 ratio of price to annual rent. Rows are ordered by price-to-rent ratio.
  \end{tablenotes}
  \end{threeparttable}
\end{table}

The population-weighted aggregate (buy-win 49\%, median $R$ 0.98) sits almost exactly on the representative-dwelling figures (51\%, 0.98). The two part company on the pound gap ($+$\gbp{1{,}263} against $-$\gbp{994}) and on appreciation (3.12\% against 3.21\%), because averaging nine regional medians and nine regional growth rates is not the same operation as measuring the composite series itself.

The cross-section tracks affordability only loosely. The West Midlands, the East Midlands and the North West (price-to-rent ratios of 20.1, 21.8 and 19.1) delivered the highest buy-win shares at 63\%, 56\% and 55\%, and the largest median gaps in the buyer's favour. The South East, on the highest ratio at 22.4, managed 42\%. But the North East breaks the pattern outright: it has the cheapest housing relative to rent of any region, at 17.7, and the \emph{lowest} buy-win share at 38\%.

Splitting cohorts by entry era shows where the North East diverges. In the post-crisis window of 2010--2015, when cheap debt made buying win almost everywhere, North East buyers won just 14\% of cohorts against 82\% in the West Midlands; from 2016 onwards both won essentially every cohort, the North East 98\% and the West Midlands all of them. That runs against the national record, where renting won every 2016 entry: the representative dwelling is weighted towards the South East and London, where those cohorts lost. North East prices simply did not recover after 2009, appreciating 2.13\% per year across the sample against 3.81\% in London. A low price-to-rent ratio makes owning cheap relative to renting, but it cannot substitute for capital growth once transaction costs have to be recovered.

The North East also shows what leverage does when appreciation fails. Eleven of its 199 cohorts, every one of them entering between June 2007 and June 2008, ended a five-year hold with \emph{negative} net worth after selling costs, and five North West cohorts, entering between June 2007 and February 2008, did the same. No other region produced a single instance. A 10\% deposit in a low-appreciation region is a position an ordinary downturn can destroy outright, and a wealth ratio conceals that in a way a pound figure does not.

\subsection{What entry conditions determine}
\label{sec:predictors}

I regress the five-year cohort wealth ratio on the entry mortgage rate and the entry price-to-rent ratio.

\begin{center}
\begin{threeparttable}
\begin{tabular}{lrrr}
\toprule
 & \textbf{Pearson $r$} & \textbf{$R^2$} & \textbf{HAC $t$} \\
\midrule
Mortgage rate       & 0.42 & 0.18 & 3.42 \\
Price-to-rent ratio & 0.03 & 0.00 & 3.04 \\
Joint (OLS)         & --   & 0.56 & --   \\
\bottomrule
\end{tabular}
\begin{tablenotes}
  \small
  \item \textit{Notes:} 199 cohorts, each holding five years. HAC $t$-statistics use Hansen--Hodrick standard errors with the truncation set to the overlap length (59 months); the corresponding OLS statistics are 15.9 and 13.2, and the difference is the cost of ignoring the overlap.
\end{tablenotes}
\end{threeparttable}
\end{center}

Jointly the two variables account for 56\% of the cross-cohort variation. Pearson $r$ and its $R^2$ above are bivariate while the $t$-statistics are on the joint regression, which is why the price-to-rent ratio shows a $t$ of 3.04 despite a bivariate $r$ of 0.03 and $p$ of 0.70: the two regressors are collinear, and the ratio contributes through the joint fit rather than on its own. Three further qualifications change what the 56\% is a claim about.

\emph{First, most of the relationship is an accounting identity rather than a forecast.} The entry price-to-rent ratio is, by construction, the initial rent the renter pays relative to the price the buyer pays; it therefore fixes the monthly cash-flow difference that the renter invests, and the entry rate fixes the buyer's payment for the first fix. To measure how much of the fit this mechanism alone produces, I re-ran all 199 cohorts in a counterfactual world with nothing whatever to forecast: house prices, rents and equities all grow at their sample-mean rates and the mortgage rate never moves, so every cohort faces an identical future and differs only in where it entered. In that world the same two variables explain \textbf{97\%} of the variation in outcomes. The regression is largely re-deriving the user-cost identity of Section~\ref{sec:rar}, not discovering a predictive relationship.

\emph{Second, a high $R^2$ is what overlapping cohorts produce even when nothing is predictable.} Because the 199 cohorts overlap and both regressors are highly persistent, the sampling distribution of $R^2$ under the null of no relationship is not centred near zero. Shifting the predictor series circularly against the outcomes, which preserves each series' own autocorrelation exactly and destroys only their alignment, gives, across 5{,}000 draws, a median $R^2$ of 0.29 and a 99th percentile of 0.58 under no predictability at all. The observed 0.56 clears that null ($p = 0.034$), but only just, and reporting 0.56 without the benchmark would overstate it substantially. The Hansen--Hodrick correction tells the same story: the $t$-statistics fall from about 15 to about 3.

\emph{Third, there is no out-of-sample forecasting power in this sample.} A buyer in real time could only have estimated the relationship on cohorts whose five years had already elapsed. Imposing that constraint, the expanding-window forecast is \emph{worse} than simply using the historical mean outcome, with a Campbell--Thompson out-of-sample $R^2$ of $-2.5$ and a Clark--West statistic of $-2.1$. What binds is the number of independent cohorts rather than the number of cohorts: 258 months contain only 4.3 non-overlapping five-year windows, and no amount of monthly restarting creates more.

Through the arithmetic Section~\ref{sec:rar} makes explicit, these two variables largely \emph{determine} the outcome without forecasting it: buy when the price-to-rent ratio is low and financing is cheap, and the cash flows work in your favour because you have bought cheaply relative to renting.

One empirical regularity does survive, and it is one-sided (Figure~\ref{fig:predictors}). Renting won \emph{every} cohort that entered at a price-to-rent ratio above 24.4. There is no corresponding floor: renting also won at the lowest ratio in the sample, 17.6, because the 2007--2008 entrants bought at moderate ratios and were still overwhelmed by the crash. A floor does appear if cohorts are evaluated at a common sample end rather than over a fixed hold, but it is an artefact of that design. The representative ratio averaged 21.9 over the sample, peaked at 26.4 in August 2022 and stood at 20.7 in June 2026, below the ceiling and in the contested middle (Figure~\ref{fig:market_overview}, Panel C).

\begin{figure}[H]
  \centering
  \includegraphics[width=\textwidth]{fig5_predictors}
  \caption{Conditions at entry and the eventual outcome, 199 cohorts each holding five years. Panel A: quoted first-time-buyer mortgage rate at entry. Panel B: price-to-rent ratio at entry. Filled markers denote cohorts in which buying produced greater terminal wealth, hollow markers those in which renting-and-investing did. Dashed lines are ordinary least-squares fits.}
  \label{fig:predictors}
\end{figure}

% ============================================================================
% 06_rar.tex — The required appreciation rate
% ============================================================================

\section{The required appreciation rate}
\label{sec:rar}

The cohort analysis of Section~\ref{sec:results} looks backwards. It tells a buyer what would have happened had they bought at some past date, which is not the question they face. This section turns the same model forwards.

Define the \emph{required appreciation rate} (RAR) as the constant annual house-price growth at which buying exactly matches renting-and-investing over a given holding period, with every other input held at its assumed path. It is the user-cost equilibrium of \textcite{poterba1984} and \textcite{himmelberg2005} solved for the one term those authors have to assume. Rather than positing an expected capital gain and pronouncing prices fair or unfair, the RAR inverts the question: it reports the growth rate buying needs, and leaves the reader to supply their own expectation. If prices grow faster than the RAR the buyer wins; slower, and the renter wins.

The RAR is a threshold, so it requires no view on house prices. The one judgement it does need is the return the renter's portfolio earns instead, and since that judgement drives the answer, this section reports it under two openly different assumptions.

\subsection{Assumptions}
\label{sec:rar_assumptions}

The calculation uses the last month of the sample, June 2026: a representative English dwelling priced at \gbp{319{,}488} let at \gbp{1{,}286} per month, a price-to-rent ratio of 20.7 and a gross rental yield of 4.83\%. First-time-buyer stamp duty on that price is \gbp{974}.

Every parameter of the baseline in Table~\ref{tab:assumptions} carries over unchanged: a 10\% deposit, a 30-year term, maintenance and depreciation at 1.5\% of value per year, purchase costs of 1.2\% and selling costs of 1.8\%. Two inputs are worth stating explicitly.

\emph{The mortgage rate} is 4.94\%, the June 2026 quoted rate on a five-year fix at 90\% loan-to-value (the same series the historical analysis uses, at its latest observation). It is held constant across the horizon. This is a real assumption and it cuts against the buyer at long horizons: a borrower refixing every five years is exposed to rates that may be higher or lower, and holding today's rate fixed for ten years understates that uncertainty in both directions. Rent is assumed to grow at 3\% per year, against the 2.9\% the representative dwelling's rent actually delivered over the sample.

\emph{The renter's return} is the input that matters most, and there is no single defensible value for it. The paper therefore reports two.

The \textbf{cautious} figure is 5.9\% nominal: a CAPE-implied long-run real return on global equities of roughly 3\%, plus expected inflation of roughly 3\%, less the 0.12\% fund fee. The \textbf{realised} figure is 10.4\%, the return the renter's portfolio actually earned over 2005--2026 net of that fee. They are two coherent views rather than bounds on a distribution, one anchored on current valuations and one on recent experience. Reporting only the first flatters the buyer; only the second flatters the renter. The 4.94\% ten-year gilt yield provides a third, lower-risk comparator.

\subsection{Results}

\begin{table}[H]
  \centering
  \caption{Forward-looking required appreciation rate, June 2026 conditions}
  \label{tab:rar_horizon}
  \begin{threeparttable}
  \begin{tabular}{@{}rrrr@{}}
    \toprule
    & \multicolumn{2}{c}{\textbf{vs.\ global equities}} & \\
    \cmidrule(lr){2-3}
    \textbf{Horizon} & \textbf{at 5.9\% (cautious)} & \textbf{at 10.4\% (realised)} & \textbf{vs.\ gilts (4.9\%)} \\
    \midrule
    3 years  & 2.77\% & 3.53\% & 2.62\% \\
    5 years  & 2.27\% & 3.24\% & 2.08\% \\
    7 years  & 2.03\% & 3.22\% & 1.80\% \\
    10 years & 1.79\% & 3.38\% & 1.50\% \\
    \bottomrule
  \end{tabular}
  \begin{tablenotes}
    \small
    \item \textit{Notes:} Constant annual house-price growth at which buying exactly matches renting-and-investing, at June 2026 inputs with 3\% annual rent growth and the mortgage rate held at 4.94\%. Baseline parameters of Table~\ref{tab:assumptions}. For comparison, the representative English dwelling appreciated at 3.21\% per year over 2005--2026.
  \end{tablenotes}
  \end{threeparttable}
\end{table}

The two equity assumptions bracket the historical outcome. Over a five-year hold, the horizon of the headline result, buying needs 2.27\% annual appreciation to match a portfolio compounding at the cautious 5.9\%, comfortably below the 3.21\% the dwelling actually delivered. Against a portfolio compounding at the realised 10.4\%, it needs 3.24\%, which is the realised rate almost exactly. Whether buying looks attractive today therefore depends almost entirely on what one expects of global equities, and hardly at all on one's view of house prices.

Figure~\ref{fig:rar_fwd} shows why the horizon matters. Under the cautious assumption the requirement falls throughout, from 2.77\% at three years to 1.79\% at ten: transaction costs amortise, the mortgage amortises, and a rising rent is compared against a payment schedule fixed in nominal terms. Under the realised assumption a portfolio compounding at 10.4\% pulls away fast enough to offset those horizon effects, and the requirement is close to flat: 3.53\% at three years, 3.24\% at five, a minimum of 3.22\% at seven, then back up to 3.38\% at ten. Across the three-to-ten-year range the cautious requirement stays below the 3.2\% the dwelling actually delivered and the realised requirement stays above it, so the horizon moves the numbers without changing which assumption favours which side.

\begin{figure}[H]
  \centering
  \includegraphics[width=\textwidth]{fig6_rar}
  \caption{Forward-looking required appreciation rate by holding period, at June 2026 conditions. Each line gives the constant annual house-price growth at which buying exactly matches renting-and-investing, against a global equity portfolio earning the CAPE-implied 5.9\% or the 10.4\% realised over 2005--2026, or against the 4.9\% ten-year gilt yield. The dotted reference is the 3.2\% per year the representative English dwelling actually appreciated over the sample. Sources: HM Land Registry UK HPI; ONS PIPR; Bank of England; MSCI.}
  \label{fig:rar_fwd}
\end{figure}

The five-year requirement of 2.27\% has a simple user-cost reading. Renting costs the gross rental yield, 4.83\% of the property's value per year. Owning costs mortgage interest on the 90\% loan (about $4.94\% \times 0.9$), the opportunity cost of the 10\% deposit, and 1.5\% maintenance, together roughly 7.0\% before any appreciation, against which the buyer sets the rent they no longer pay. The residual is the growth buying needs, and the horizon effects above account for its decline.

\subsection{Regional break-evens}

The same calculation region by region, with each region's June 2026 price and rent, the national mortgage rate and region-specific stamp duty, gives a five-year requirement ranging from 1.17\% in the North East to 2.90\% in the South East under the cautious assumption, and from 2.04\% to 3.94\% under the realised one. Figure~\ref{fig:rar_map} maps the gradient.

Comparing each region's break-even with its own realised appreciation is where the choice of equity assumption does its work. Under the cautious assumption every one of the nine regions' break-evens sits below what it actually delivered, by margins of 0.2 to 1.6 percentage points. Under the realised assumption the ordering reverses in six: Yorkshire and the Humber, the South West, the East Midlands, the East of England, London and the South East would all have needed more appreciation than they produced. Only the three regions with the lowest price-to-rent ratios (the North East, the North West and the West Midlands) clear the bar on both assumptions.

The gradient tracks the price-to-rent ratio rather than the price level, with one exception. London has the third-lowest ratio of the nine at 20.1, but the second-\emph{highest} break-even at 2.82\%, because at \gbp{553{,}870} the average London dwelling sits above the \gbp{500{,}000} ceiling at which first-time-buyer stamp-duty relief is withdrawn altogether. That single threshold costs a London first-time buyer \gbp{17{,}694} against nothing at all in five of the nine regions, which moves the region three places up the ranking. What makes London expensive to enter is the stamp duty rather than the valuation.

\begin{figure}[H]
  \centering
  \includegraphics[width=0.78\textwidth]{fig8_regional_map}
  \caption{Forward-looking five-year required appreciation rate against global equities at the cautious 5.9\% assumption, by English region, at June 2026 conditions (per cent per year). Darker shading indicates a higher break-even, that is, a market in which buying is harder to justify. Values are those reported in Table~\ref{tab:regional}.}
  \label{fig:rar_map}
\end{figure}

% ============================================================================
% 07_robustness.tex — Robustness
% ============================================================================

\section{Robustness}
\label{sec:robustness}

Every value in the baseline of Table~\ref{tab:assumptions} is typical rather than universal. Tables~\ref{tab:sensitivity_main} and~\ref{tab:sensitivity_other} vary each parameter in turn, reporting the January 2005 full-period ratio, the median ratio across the 199 fixed five-year cohorts, and the share of those cohorts won by buying. The first column is a single cohort held twenty-one years, an illustration on the same footing as Section~\ref{sec:single} rather than a result at that horizon. The tables are generated directly from the simulation output rather than transcribed.

\begin{table}[H]
  \centering
  \caption{Sensitivity to maintenance and deposit}
  \label{tab:sensitivity_main}
  \begin{threeparttable}
  \begin{tabular}{@{}lrrr@{}}
    \toprule
    \textbf{Parameter} & \textbf{$R$ (2005)} & \textbf{Median $R$} & \textbf{Buy-win \%} \\
    \midrule
    \multicolumn{4}{@{}l}{\textit{Maintenance and depreciation, UK-grounded range (\% of value p.a.)}} \\
    \quad 1.2 & 1.24 & 0.92 & 63 \\
    \quad 1.5 (baseline) & 1.44 & 0.98 & 51 \\
    \quad 1.8 & 1.65 & 1.04 & 46 \\
    \addlinespace
    \multicolumn{4}{@{}l}{\textit{Maintenance and depreciation, beyond the UK evidence}} \\
    \quad 1.0 & 1.10 & 0.87 & 71 \\
    \quad 2.0 & 1.79 & 1.08 & 44 \\
    \quad 2.5 & 2.14 & 1.19 & 36 \\
    \quad 3.0 & 2.48 & 1.28 & 18 \\
    \addlinespace
    \multicolumn{4}{@{}l}{\textit{Deposit (\% of purchase price)}} \\
    \quad 5 & 1.42 & 1.03 & 46 \\
    \quad 10 (baseline) & 1.44 & 0.98 & 51 \\
    \quad 15 & 1.52 & 0.98 & 54 \\
    \quad 20 & 1.59 & 0.97 & 54 \\
    \quad 25 & 1.66 & 0.99 & 52 \\
    \quad 50 & 2.00 & 1.11 & 20 \\
    \quad 100 (cash) & 2.59 & 1.22 & 8 \\
    \bottomrule
  \end{tabular}
  \begin{tablenotes}
    \small
    \item \textit{Notes:} $R > 1$: renting-and-investing produced greater terminal wealth. $R$ (2005) is the January 2005 entry cohort held to June 2026. Median $R$ and buy-win \% are taken over the 199 fixed five-year cohorts. All other parameters at baseline. The maintenance range of 1.2--1.8\% is what UK national-accounts evidence supports (Section~\ref{sec:framework}); the values below and above it are reported as bounds, not as equally supported alternatives.
  \end{tablenotes}
  \end{threeparttable}
\end{table}

\begin{table}[H]
  \centering
  \caption{Sensitivity to other parameters}
  \label{tab:sensitivity_other}
  \begin{threeparttable}
  \begin{tabular}{@{}lrrr@{}}
    \toprule
    \textbf{Parameter} & \textbf{$R$ (2005)} & \textbf{Median $R$} & \textbf{Buy-win \%} \\
    \midrule
    \multicolumn{4}{@{}l}{\textit{Mortgage term}} \\
    \quad 25 years & 1.56 & 1.00 & 51 \\
    \quad 30 years (baseline) & 1.44 & 0.98 & 51 \\
    \quad 35 years & 1.36 & 0.97 & 53 \\
    \addlinespace
    \multicolumn{4}{@{}l}{\textit{Rate fixation}} \\
    \quad 2-year fixes & 1.40 & 0.87 & 65 \\
    \quad 5-year fixes (baseline) & 1.44 & 0.98 & 51 \\
    \quad 10-year fixes & 1.46 & 0.98 & 51 \\
    \addlinespace
    \multicolumn{4}{@{}l}{\textit{Mortgage rate series}} \\
    \quad FTB 5y fix at 90\% LTV (baseline) & 1.44 & 0.98 & 51 \\
    \quad CFMBJ95 effective rate, floating & 1.20 & 0.74 & 74 \\
    \addlinespace
    \multicolumn{4}{@{}l}{\textit{Renter's leverage}} \\
    \quad 1.00$\times$ (baseline, unlevered) & 1.44 & 0.98 & 51 \\
    \quad 1.25$\times$ & 1.57 & 1.02 & 48 \\
    \quad 1.50$\times$ & 1.69 & 1.07 & 45 \\
    \quad 1.75$\times$ & 1.81 & 1.11 & 42 \\
    \quad 2.00$\times$ & 1.93 & 1.15 & 38 \\
    \addlinespace
    \multicolumn{4}{@{}l}{\textit{Transaction costs}} \\
    \quad Purchase costs 1.0\% & 1.43 & 0.97 & 54 \\
    \quad Purchase costs 1.2\% (baseline) & 1.44 & 0.98 & 51 \\
    \quad Purchase costs 2.0\% & 1.49 & 1.02 & 49 \\
    \quad Selling costs 1.0\% & 1.43 & 0.96 & 57 \\
    \quad Selling costs 1.8\% (baseline) & 1.44 & 0.98 & 51 \\
    \quad Selling costs 4.0\% & 1.49 & 1.07 & 46 \\
    \quad Stamp duty excluded & 1.39 & 0.95 & 57 \\
    \addlinespace
    \multicolumn{4}{@{}l}{\textit{Renter's portfolio}} \\
    \quad MSCI ACWI net, less fee (baseline) & 1.44 & 0.98 & 51 \\
    \quad MSCI ACWI gross, no fee & 1.63 & 1.01 & 49 \\
    \quad FTSE 100 total return & 0.77 & 0.76 & 64 \\
    \bottomrule
  \end{tabular}
  \begin{tablenotes}
    \small
    \item \textit{Notes:} As Table~\ref{tab:sensitivity_main}. The renter-leverage rows charge margin interest at Bank Rate plus 1.50\% and impose a 30\% maintenance margin with forced liquidation; no cohort is liquidated, for the reason given in the text.
  \end{tablenotes}
  \end{threeparttable}
\end{table}


Five regularities organise them, and they are ordered by how much they matter.

\subsection{Maintenance moves the answer most}

Across the range UK national-accounts evidence supports, 1.2\% to 1.8\% of value per year, the buy-win share moves from 63\% to 46\% and the median ratio from 0.92 to 1.04. The verdict shifts from ``buying modestly ahead'' to ``renting modestly ahead'' inside the defensible band, so the near-parity headline is conditional on this calibration.

Two points bound how much weight that deserves. The supported band does not run from 1\% to 3\%: the ONS depreciation figure is 1.20\% of market value on average and 1.47\% at its 1995 peak, and adding cash maintenance and insurance cannot reach 2\% without double-counting capital improvements as running costs. Within the band the buy-win share moves 17 percentage points and crosses 50\%, so the sign of the verdict is not safe inside the supported range. Across the unsupported 1\%--3\% range it would move from 71\% to 18\%.

\subsection{Leverage helps the buyer up to about 15 per cent down}

An unlevered buyer is comprehensively beaten: at a 100\% deposit, buying wins 8\% of cohorts and the January 2005 ratio is 2.59. At a 50\% deposit it wins 20\%. Leverage is what allows a claim on an asset appreciating at 3.2\% per year to compete with a portfolio compounding at 10.4\%, and the result echoes the portfolio-choice observation that owner-occupied housing is attractive largely through its financing \parencite{flavin2002}.

But the relationship is not monotonic. Buying does \emph{best} at deposits of 15--20\% (54\% of cohorts), slightly worse at the 10\% baseline (51\%), and worse again at 5\% (46\%). The reason is visible in Section~\ref{sec:rolling}: leverage magnifies losses more than gains, and the cohorts that bought into the 2007--2008 peak were destroyed in proportion to how levered they were. A 5\% deposit produces both the best and the worst outcomes, and the median cohort is worse off for it. The first-time buyer's high loan-to-value both makes ownership competitive and makes badly-timed ownership catastrophic.

\subsection{Fixing the rate cost buyers 23 percentage points}

Replacing the first-time buyer's five-year fix with the Bank of England's effective rate on all new advances, recast monthly, moves the buy-win share from 51\% to \textbf{74\%}, a larger effect than any parameter except maintenance. Shortening the fix to two years moves it to 65\%.

That gap is the cost of fixing. A borrower who fixed for five years in 2005 paid a pre-crisis rate through to 2010 and did not receive the 2009 collapse; a floating borrower received it within months. Over this particular sample, in which rates fell for a decade, fixing was expensive. That is a statement about 2005--2026 rather than a recommendation: the same design would have protected a borrower through 2022--2023, and the sample contains only one such episode. The direction of the effect is worth keeping in view when reading any buy-versus-rent comparison built on an all-borrower average rate, which will flatter buyers relative to the product a first-time buyer actually holds.

\subsection{A home-biased renter loses}

Substituting the FTSE~100 for the global portfolio moves the buy-win share from 51\% to 64\% and the January 2005 ratio from 1.44 to 0.77, turning a renter who finished 44\% ahead into one who finished 23\% behind. The three-point annual gap between UK large-cap and global equity returns (7.1\% against 10.4\%) compounds into the difference between the two verdicts. The renter's side of the parity is earned by disciplined, globally diversified investing, not by renting as such; a home-biased renter loses.

Letting the renter borrow to invest works in the same direction, as expected: at 2.00$\times$ leverage, charging margin interest at Bank Rate plus 1.50\%, the buy-win share falls to 38\%. Two caveats belong with that figure. No cohort is force-liquidated even at 2$\times$, but that is a property of \emph{this} design rather than of leverage. Because the renter is cash-flow matched, the contributions made over sixty months exceed the entire initial stake: \gbp{32{,}764} against \gbp{24{,}314} for the November 2007 cohort, the worst entry point in the sample. The position therefore deleverages itself well before a margin call can bite. The minimum equity ratio across all cohorts is 31.9\%, uncomfortably close to the 30\% maintenance margin, and monthly observation understates drawdowns: month-end missed the daily trough by roughly four percentage points in the financial crisis and nine to eleven in 2020, while brokers liquidate in real time. A daily simulation would very likely show forced sales. Separately, 2$\times$ survived 2020 only because sterling fell about 13\% in a global equity crash, which is a property of the currency pair in that episode and not of equity leverage.

\subsection{Transaction costs, stamp duty and the term are second-order}

None of these moves the buy-win share by more than a handful of percentage points across ranges wider than any real buyer faces. Doubling purchase costs from 1.0\% to 2.0\% costs five points; varying selling costs from 1.0\% to 4.0\% costs eleven; removing stamp duty entirely gains six. England's low round-trip transaction costs are a structural advantage of ownership, and the period-accurate stamp-duty module, while conceptually important, changes little at average first-time-buyer prices, though not in London, where the relief cliff above \gbp{500{,}000} is worth \gbp{17{,}694} (Section~\ref{sec:rar}). The mortgage term is almost inert: 25, 30 and 35 years give 51\%, 51\% and 53\%.

\subsection{First-time-buyer prices, matched to a first-time-buyer rent}

The headline uses the all-dwellings price against the all-property rent. Both sides describe the same dwelling, but not specifically the dwelling a first-time buyer occupies. Substituting the UK HPI first-time-buyer price series tests that.

It cannot be done by swapping the price alone. The first-time-buyer price is 84\% of the all-dwellings price, so pairing it with a whole-market rent prices a cheap dwelling against an expensive tenancy, the same composition error Section~\ref{sec:representative} identifies, and a large one: it puts the price-to-rent ratio at 17.4 against a regional range of 17.7--22.4 and makes buying win 95\% of cohorts. The rent must be matched. Applying the same principle as before, the one-bedroom PIPR rent is 82\% of the all-property rent, which tracks the 84\% price discount closely; two-bedroom, at 92\%, does not.

On that matched pair, buying wins \textbf{49\%} of cohorts against the headline's 51\%, with a median ratio of 1.02 against 0.98. The conclusion is unchanged. On the two-bedroom basis buying wins 80\%, which brackets the answer from the other side and shows that the bedroom choice matters more than the price series does.

The check runs only from January 2015, because the first-time-buyer price series begins in 2012 and the PIPR bedroom series in 2015. It therefore omits the financial crisis and the 2010--2014 recovery, the episodes that generate most of the dispersion in the headline, which is why it is a robustness check and not the main specification.

\subsection{Cohorts evaluated at the sample end}

Running every cohort forward to the last month of the sample, rather than for a fixed period, gives 255 cohorts and a buy-win share of 16\% against 51\% over fixed five-year holds. It is the design most buy-versus-rent backtests use, so the gap is worth explaining.

The gap is not a finding about tenure. Running to a common end date produces cohorts of wildly unequal length, from 258 months down to four, and the recent, short ones cannot amortise their round-trip transaction costs, which consume roughly thirty per cent of the initial outlay. A buyer who entered in 2025 and is measured eighteen months later has lost on arithmetic that has nothing to do with the merits of owning. Because entry conditions are correlated with cohort length under that design, it also contaminates the regression of Section~\ref{sec:predictors}: cohort length alone accounts for a third of the variance in outcomes. The fixed-horizon design avoids both problems.

% ============================================================================
% 08_discussion.tex — Discussion
% ============================================================================

\section{Discussion}
\label{sec:discussion}

\subsection{International context}

The near-parity result is a product of the particular history England delivered. Table~\ref{tab:bis} places English house-price growth in international context using BIS residential property price statistics \parencite{bis_rpp}. The United Kingdom sits mid-table over 2005--2025: cumulative nominal growth of 95\% (3.3\% per year), far below the central European booms and well below the pace that would have made ownership financially dominant. 

\begin{table}[H]
  \centering
  \caption{Nominal house-price growth, selected countries, 2005Q1--2025Q4}
  \label{tab:bis}
  \begin{threeparttable}
  \begin{tabularx}{0.7\textwidth}{@{}Xrr@{}}
    \toprule
    \textbf{Country} & \textbf{Cumulative} & \textbf{Annualised} \\
    \midrule
    Hungary          & 356\% & 7.6\% \\
    Portugal         & 156\% & 4.6\% \\
    Netherlands      & 110\% & 3.6\% \\
    United Kingdom   & 95\%  & 3.3\% \\
    United States    & 89\%  & 3.1\% \\
    Germany          & 85\%  & 3.0\% \\
    Ireland          & 62\%  & 2.3\% \\
    France           & 60\%  & 2.3\% \\
    Spain            & 60\%  & 2.3\% \\
    Italy            & 15\%  & 0.7\% \\
    \bottomrule
  \end{tabularx}
  \begin{tablenotes}
    \small
    \item \textit{Source:} BIS residential property prices, nominal, selected series. The UK entry covers the whole United Kingdom; the representative English series of this paper grew 96\% (3.2\% per year) over the same span.
  \end{tablenotes}
  \end{threeparttable}
\end{table}

\subsection{Limitations}

Several limitations qualify the results. \emph{Aggregation.} The representative dwelling is a construct; households decide in a specific market, and the regional results of Section~\ref{sec:regional}, with buy-win shares from 38\% to 63\%, may differ from the aggregate result. Even within regions, composition differences between the owned and rented stocks mean the price-to-rent \emph{level} carries measurement error. \emph{Behavioural assumptions.} The renter invests 100\% of every monthly difference, immediately and for two decades, in a globally diversified portfolio. The saving and portfolio-behaviour literature \parencite{laibson1997,thaler2004} puts this well above typical conduct; the FTSE-only robustness check shows how costly even a diversification failure is, before any failure of discipline. \emph{Taxes.} The symmetric no-tax treatment is appropriate for a typical first-time buyer (no capital-gains tax on a main residence; ISA capacity covering the renter's flows for most of the sample) but understates the drag on very large taxable portfolios, which would shade outcomes towards buying. \emph{Inference.} The 199 cohorts overlap and share the sample's single realisation of asset returns: 258 months contain only 4.3 non-overlapping five-year windows, so cohort counts are counts of simulations rather than of independent observations. Win shares and correlations describe this history and are not tests on independent draws. \emph{Unpriced attributes.} Ownership insures against rent risk \parencite{sinai2005} and carries control rights; renting offers more flexibility and mobility \parencite{ehs2024}. None of these enters the terminal-wealth accounting. 

\subsection{Implications}

Three implications follow for households and for the framing of housing-policy debates. First, the popular certainty that buying is the financially correct choice is unwarranted, and not only at the five-year horizon. Across the 199 five-year cohorts the two strategies finish level: a median wealth ratio of 0.98, a median gap of \gbp{994} in the owner's favour and a mean gap of \gbp{1{,}744}. Lengthening the hold tilts the count towards buying, which wins 62\% of seven-year and 61\% of ten-year cohorts, but the median renter still finishes within 8--11\% of the owner, a gap of \gbp{14{,}454} at ten years on a dwelling worth \gbp{319{,}488}. Buying is ahead more often than not beyond five years, and it is not ahead by much.

Second, entry valuation and financing cost largely determine the cohort outcome without forecasting it, for the reasons Section~\ref{sec:predictors} sets out. The one regularity that survives runs in a single direction: renting won every cohort entering above a price-to-rent ratio of 24.4, and there is no matching floor.

Third, the comparison is conditional on conduct: the renter's side of the parity requires actually investing the cost difference, every month. A mortgage, besides providing crucial leverage, functions as a commitment device, since amortisation is forced saving. For households unlikely to match that discipline voluntarily, home ownership is likely the safer default even where the arithmetic is level, and that behavioural asymmetry is probably the strongest practical argument for buying.
% ============================================================================
% 09_conclusion.tex — Conclusion
% ============================================================================

\section{Conclusion}
\label{sec:conclusion}

This paper provides a first-time-buyer financial comparison of buying and renting-with-investing on English data, run over every monthly entry point from January 2005 to July 2021 with a five-year hold, on a representative dwelling whose price and rent share one regional composition. The answer to the popular question is that, over these two decades, there was close to no answer: buying produced greater terminal wealth in 102 of 199 cohorts (51\%), the median wealth ratio was 0.98, and the median gap \gbp{994} in the buyer's favour. Entry conditions (valuation and financing cost at the point of purchase) decided which side of parity a household landed on: renting won every cohort entering between 2006 and 2009 and again in 2016, buying won every cohort entering from 2011 to 2013 and in 2018 and 2019, and the regional buy-win share ran from 38\% in the North East to 63\% in the West Midlands.

That comparability extends well past five years. Buying won 62\% of seven-year holds and 61\% of ten-year holds, and the median renter still finished within 8--11\% of the owner, \gbp{14{,}454} behind at ten years on a dwelling worth \gbp{319{,}488}. Beyond five years buying is ahead more often than not, by margins small enough that renting-and-investing remains a defensible financial choice for a household that keeps to it. The distribution around that parity is asymmetric: buying won more cohorts, renting won the larger amounts. Renting more than doubled the owner's wealth in 21 of the 199 five-year cohorts and more than quadrupled it in five, all of which entered between August and December 2007, peaking at a ratio of 4.17 for the November 2007 cohort. That tail is why the mean ratio of 1.22 favours the renter while the mean gap of \gbp{1{,}744} favours the owner: renting's proportional wins landed where the owner's equity had collapsed to a small number, and buying's wins landed on much larger balance sheets.

At the short end the arithmetic is a judgement on transaction costs rather than on ownership. They consume roughly thirty per cent of a first-time buyer's initial outlay on the round trip and take about six years to amortise. Buying won only 30\% of one-year holds, so a household expecting to move within two or three years is, on this evidence, usually better off renting. The shape of that profile belongs to this sample, which contains one crash, one long expansion of cheap credit and one exceptional equity bull market, and I would not read it as a general fact about tenure.

Beyond the headline, the analysis yields a methodological caution and a forward-looking metric. The caution is that the official national house-price and rent aggregates weight regions differently, so their ratio (the most widely quoted price-to-rent figure) lies outside the range of every English region and flatters ownership in any cost comparison built on it. The metric is the required appreciation rate. At June 2026 conditions a five-year buyer needs 2.3\% annual appreciation to match a cautiously parameterised global equity portfolio and 3.2\% to match the return that portfolio actually delivered, while the dwelling itself appreciated at 3.2\%. The forward-looking arithmetic is therefore genuinely balanced, in the specific sense that plausible assumptions fall on both sides of it. Reporting only the cautious figure would make the case for buying look settled when it is not.

The results also discipline how the trade-off should be described. Leverage is what keeps buying competitive, rather than house-price growth: an unlevered cash buyer wins only 8\% of cohorts. But more leverage is not simply better. Buying does best at deposits of 15--20\% and worse at 5\%, because leverage magnifies the losses of a badly-timed purchase more than it magnifies the gains of a well-timed one. In the North East and North West, sixteen cohorts entering between mid-2007 and mid-2008 ended their five years with negative net worth after selling costs. Conversely, the renter's parity is earned by conduct: it requires investing the full cost difference in a \emph{globally} diversified portfolio, and a renter who bought only UK equities loses. Two calibrations move the answer more than any others: the maintenance-and-depreciation charge and the fixation design. Section~\ref{sec:robustness} reports the full mapping from each to the headline so that readers can relocate the result under their own priors.

All data, code, and an interactive implementation of the framework are openly available, and the analysis is designed to be re-run as new months of official data arrive. Extensions to Scotland and Wales as their rent statistics mature, first-time-buyer-specific price series matched to an appropriate rent, and the related buy-to-let decision are natural next steps.


% --- Bibliography ---
\newpage
\printbibliography[title={References}]

% --- Appendix ---
\newpage
% ============================================================================
% appendix.tex — Appendices
% ============================================================================

\begin{appendices}

\section{Data construction}
\label{app:data}

\subsection*{Rents}

The rent series splices two ONS products. The PIPR historical series (released March 2025) provides average monthly rents in GBP for England and the nine English regions from January 2005 to February 2025, chain-linked to the predecessor Index of Private Housing Rental Prices; the live PIPR monthly dataset provides the same measures from January 2015 onwards. Over their 122-month overlap (January 2015--February 2025) the two agree exactly at every observation, so the splice, historical values before January 2015 and live values thereafter, introduces no discontinuity. Rent levels are model-consistent estimates published by the ONS (base-period mean rents uprated by the index), not raw sample means.

\subsection*{House prices}

Regional and national average prices (all dwellings) are taken directly from the UK HPI full file (June 2026 vintage). The UK HPI revises its most recent months; the June 2026 observation used here is the first published estimate.

\subsection*{Representative dwelling}

The representative price and rent are fixed-weight averages of the nine regional series with mid-2023 ONS population weights (millions): North East 2.65, North West 7.52, Yorkshire and The Humber 5.56, East Midlands 5.02, West Midlands 6.11, East of England 6.40, London 8.87, South East 9.46, South West 5.80. Using private-rented-sector household counts instead of population moves the June 2026 representative price-to-rent ratio from 20.7 to 20.6, and no headline result materially.

\subsection*{Mortgage rate}

The baseline rate is a five-year fix at 90\% loan-to-value, refixed every sixty months and repriced at each refix by the borrower's current loan-to-value (Section~\ref{sec:mortgage_rate}). No single Bank of England series covers it, so it is constructed from the quoted-rate family \parencite{boe_quoted}, using five-year data wherever possible so that no assumption about the term structure is needed:

\begin{itemize}
  \item \textbf{February 2019 onwards:} the published five-year 90\% LTV series (IUMZO28) directly.
  \item \textbf{Before February 2019:} interpolated in loan-to-value between the published five-year series at 75\% (IUMBV42) and 95\% (IUM5WTL), weight 0.55. The weight is fitted against IUMZO28 over the ninety months where it exists; mean discrepancy 0.01 percentage points, standard deviation 0.24. It is below the 0.75 that linear interpolation implies because the premium is convex in loan-to-value.
  \item \textbf{October 2008 to October 2013:} IUM5WTL is unpublished for all sixty-one months, because fewer than three lenders offered a 95\% LTV five-year fix and the Bank suppresses averages below that threshold. The window is filled from the two-year 90\% series (IUMB482) plus 0.30 percentage points, calibrated against contemporaneous quoted products; it puts January 2010 at 6.91\% against a market range of 6.49\% to 7.49\%. IUMB482 is itself unpublished for March to May 2009 and is interpolated across those three months. The patch is not level-shifted to meet the interpolation, so there is a step of about 0.3 percentage points at each seam.
  \item \textbf{Bands.} The 75\% and 95\% series are carried through so the engine can reprice by current loan-to-value, piecewise-linearly between bands, with the 75\% rate used unchanged below 75\%.
\end{itemize}

The resulting series averages 4.74\% and ranges from 2.33\% to 6.99\%, sitting 1.19 percentage points above CFMBJ95.

Five-year data is used in preference to transplanting a two-year loan-to-value premium onto a five-year rate, because the loan-to-value premium compresses at longer tenors when high-loan-to-value credit is rationed: the Bank's data give a five-year-minus-two-year term premium of 0.45 percentage points at 75\% loan-to-value against 0.27 at 95\%. A transplant is unbiased after 2019 but puts January 2010 at 8.20\%, above anything then available.

CFMBJ95 (Bankstats table G1.4, the weighted-average rate on new advances to households, January 2004 onwards) is retained for the floating-rate case in Section~\ref{sec:robustness}. Bank Rate (IUMABEDR) supplies the margin cost in the renter-leverage case.

\subsection*{Equity returns}

MSCI publishes end-of-month ACWI index levels directly in GBP. I use the net-total-return variant (dividends reinvested after withholding tax) and convert levels to simple monthly returns. The baseline series deducts a fund fee of 0.12\% per year, applied as a constant monthly drag: $r^{\text{net}}_t = (1 + r_t)(1 - 0.0012)^{1/12} - 1$. The FTSE 100 total-return comparator is the net-asset-value total return of a low-cost FTSE 100 ETF (dividends reinvested at NAV; ongoing charge 0.40\% per year to June 2014, 0.07\% thereafter), which tracks the FTSE 100 total-return index with correlation 0.9999 over the months both are available.

\subsection*{Replication}

All raw data under the Open Government Licence, all cleaning and analysis code, and scripts that re-download the licence-restricted financial series from their public endpoints are available at \url{https://github.com/Martin-Cimprich/buy-vs-rent-england}.

\section{Stamp duty schedules}
\label{app:sdlt}

Table~\ref{tab:sdlt} summarises the SDLT liability of a first-time buyer in England by purchase month, as implemented in the simulation. Before 4 December 2014 SDLT was a ``slab'' tax (the rate applied to the whole price); thereafter a ``slice'' tax (marginal bands). Where first-time-buyer relief exists, the relief schedule applies below its price cap and the standard schedule above it.

\begin{table}[H]
  \centering
  \caption{SDLT for an England first-time buyer, by period}
  \label{tab:sdlt}
  \begin{threeparttable}
  \small
  \begin{tabularx}{\textwidth}{@{}lX@{}}
    \toprule
    \textbf{Period} & \textbf{Schedule applying to a first-time buyer} \\
    \midrule
    to 16 Mar 2005        & Slab: 0\% to \gbp{60}k; 1\% to \gbp{250}k; 3\% to \gbp{500}k; 4\% above \\
    17 Mar 2005--22 Mar 2006 & Slab, nil-rate band \gbp{120}k \\
    23 Mar 2006--2 Sep 2008  & Slab, nil-rate band \gbp{125}k \\
    3 Sep 2008--31 Dec 2009  & Slab, nil-rate band \gbp{175}k (holiday) \\
    1 Jan 2010--24 Mar 2010  & Slab, nil-rate band \gbp{125}k \\
    25 Mar 2010--24 Mar 2012 & First-time-buyer relief: 0\% to \gbp{250}k; standard slab above \\
    25 Mar 2012--3 Dec 2014  & Slab, nil-rate band \gbp{125}k (relief withdrawn; 7\% top band from Mar 2012) \\
    4 Dec 2014--21 Nov 2017  & Slice: 0\% to \gbp{125}k; 2\% to \gbp{250}k; 5\% to \gbp{925}k; 10\% to \gbp{1.5}m; 12\% above \\
    22 Nov 2017--7 Jul 2020  & FTB relief: 0\% to \gbp{300}k, 5\% \gbp{300}--500k (no relief above \gbp{500}k) \\
    8 Jul 2020--30 Jun 2021  & Holiday: 0\% to \gbp{500}k for all buyers; slice above \\
    1 Jul 2021--30 Sep 2021  & Standard nil-rate \gbp{250}k; FTB relief 0\% to \gbp{300}k, 5\% to \gbp{500}k \\
    1 Oct 2021--22 Sep 2022  & Standard slice (\gbp{125}k nil); FTB relief 0\% to \gbp{300}k, 5\% to \gbp{500}k \\
    23 Sep 2022--31 Mar 2025 & FTB relief: 0\% to \gbp{425}k, 5\% \gbp{425}--625k (no relief above \gbp{625}k) \\
    1 Apr 2025--             & FTB relief: 0\% to \gbp{300}k, 5\% \gbp{300}--500k (no relief above \gbp{500}k) \\
    \bottomrule
  \end{tabularx}
  \begin{tablenotes}
    \small
    \item \textit{Sources:} HMRC SDLT rate tables 2003--2026; House of Commons Library briefings SN07050 and CBP-9814. The simulation evaluates the schedule at each month's representative (or regional) average price. Scotland (LBTT, from April 2015) and Wales (LTT, from April 2018) operate different schedules and are outside the scope of this paper.
  \end{tablenotes}
  \end{threeparttable}
\end{table}

\end{appendices}


\end{document}
